\part{Über die Zeit}

\chapter{Rückwärts denken}

\section{Der Fall}

\begin{merksatz}
\boxtitel{Welche Therapie zuerst?}
Zwei Behandlungen stehen zur Verfügung; scheitert die erste, kann man die zweite versuchen.

\begin{center}
{\small
\renewcommand{\arraystretch}{1.5}
\begin{tabular}{@{}P{2.4cm}P{2.6cm}P{2.8cm}@{}}
\toprule
\textbf{Reihenfolge} & \textbf{erste wirkt} & \textbf{zweite danach}\\
\midrule
$A$, dann $B$ & $0{,}60$ & $0{,}50$\\
$B$, dann $A$ & $0{,}40$ & $0{,}30$\\
\bottomrule
\end{tabular}}
\end{center}
\kursiv{Erfolg in der ersten Runde ist $1{,}00$ wert, in der zweiten $0{,}90$, kein Erfolg $0{,}30$.}

\begin{center}
\begin{minipage}{0.88\textwidth}
\kursiv{\textbf{Man beachte die zweite Zeile: $A$ wirkt nach Versagen von $B$ schlechter als sonst} --- etwa weil Zeit verstrichen ist oder weil $B$ die Empfindlichkeit verändert hat. Solche Abhängigkeiten sind der Grund, warum man Reihenfolgen nicht einzeln bewerten kann.}
\end{minipage}
\end{center}
\end{merksatz}

\begin{merksatz}
\boxtitel{Die Auswertung}
\schritt{Zuerst $A$}
\[ 0{,}60\cdot1{,}00 + 0{,}40\cdot\bigl[0{,}50\cdot0{,}90 + 0{,}50\cdot0{,}30\bigr] = 0{,}600 + 0{,}40\cdot0{,}600 = \mathbf{0{,}8400} \]

\schritt{Zuerst $B$}
\[ 0{,}40\cdot1{,}00 + 0{,}60\cdot\bigl[0{,}30\cdot0{,}90 + 0{,}70\cdot0{,}30\bigr] = 0{,}400 + 0{,}60\cdot0{,}480 = \mathbf{0{,}6880} \]

\begin{center}
\begin{minipage}{0.9\textwidth}
\textbf{Man rechnet von hinten: erst den Wert der zweiten Runde, dann den der ersten.}
\end{minipage}
\end{center}
\end{merksatz}

\section{Das Verfahren}

\begin{merksatz}
\boxtitel{Satz 7.1 --- Bellmans Gleichung}
Für ein Problem mit Zuständen $z$, Handlungen $a$, Übergängen $\Pr[z' \mid z,a]$ und Nutzen $u(z,a)$ über $T$ Schritte gilt
\[ V_t(z) = \max_a \left\{ u(z,a) + \sum_{z'} \Pr[z'\mid z,a]\ V_{t+1}(z') \right\}, \qquad V_T(z) = u_{\text{Ende}}(z) . \]

\schritt{Begründung}
\kursiv{Im letzten Schritt} ist $V_T$ gegeben.

\kursiv{Kennt man $V_{t+1}$}, so ist der beste Zug in $t$ derjenige, der die Summe aus sofortigem Nutzen und erwartetem künftigem Wert maximiert --- denn was nach $t$ geschieht, hängt nur vom erreichten Zustand ab.

\kursiv{Also} bestimmt man $V_{T-1}, V_{T-2}, \ldots$ der Reihe nach. $\blacksquare$

\begin{center}
\begin{minipage}{0.9\textwidth}
\textbf{Das Optimalitätsprinzip: Jede Fortsetzung einer besten Strategie ist selbst eine beste Strategie --- vom erreichten Zustand aus.}
\end{minipage}
\end{center}

\schritt{Was das erspart}
\begin{center}
{\small
\renewcommand{\arraystretch}{1.48}
\begin{tabular}{@{}P{2.6cm}P{2.6cm}P{2.6cm}@{}}
\toprule
\textbf{Schritte} & \textbf{alle Strategien} & \textbf{rückwärts}\\
\midrule
$5$ & $2^{31}$ & $5\cdot|Z|\cdot|A|$\\
$20$ & $> 10^{300}$ & $20\cdot|Z|\cdot|A|$\\
\bottomrule
\end{tabular}}
\end{center}

\begin{center}
\begin{minipage}{0.88\textwidth}
\kursiv{Eine Strategie ordnet jedem Zustand und jedem Zeitpunkt eine Handlung zu; ihre Zahl wächst doppelt exponentiell. \textbf{Die Rückwärtsrechnung wächst linear in der Zeit} --- dieselbe Ersparnis wie bei der dynamischen Programmierung in Band 19, Kapitel~11, und im Grunde derselbe Satz.}
\end{minipage}
\end{center}
\end{merksatz}

\section{Der Bruchfall}

\begin{warnung}
\boxtitelr{Wo das Prinzip versagt}
\begin{center}
{\small
\renewcommand{\arraystretch}{1.52}
\begin{tabular}{@{}P{2.8cm}P{5.8cm}@{}}
\toprule
\textbf{Kein Markow-Zustand} & Hängt der Nutzen vom ganzen Weg ab, so genügt der Zustand nicht --- man muss die Geschichte in ihn aufnehmen.\\
\textbf{Unvollständige Sicht} & Kennt man den Zustand nur ungenau, wird der \kursiv{Glaube} über ihn zum Zustand --- und der ist stetig.\\
\textbf{Nichtlineare Ziele} & Wer nicht den Erwartungswert maximiert, kann nicht rückwärts zerlegen.\\
\bottomrule
\end{tabular}}
\end{center}

\begin{center}
\begin{minipage}{0.88\textwidth}
\kursiv{Die dritte Zeile trifft manche Größen, die in der Medizin gern verwendet werden. \textbf{Wer etwa die Wahrscheinlichkeit maximieren will, mindestens fünf Jahre zu überleben, kann das rückwärts rechnen} --- das ist wieder ein Erwartungswert. Wer aber den Median der Überlebenszeit maximieren will, kann es nicht: Mediane lassen sich nicht schrittweise zerlegen.}
\end{minipage}
\end{center}
\end{warnung}

\section{Übungsaufgaben}

\begin{aufgabe}
\boxtitel{Aufgabe 7.1 --- Die Reihenfolge kippen}
Ab welcher Wirksamkeit von $B$ in der ersten Runde wäre es besser, mit $B$ zu beginnen?

\loesung
\schritt{Schritt 1 --- Der Ansatz}
\[ q\cdot1{,}00 + (1-q)\bigl[0{,}30\cdot0{,}90+0{,}70\cdot0{,}30\bigr] = 0{,}8400 \]

\schritt{Schritt 2 --- Auflösen}
\[ q + (1-q)\cdot0{,}480 = 0{,}840 \]
\[ 0{,}520\,q = 0{,}360 \quad\Longrightarrow\quad q = 0{,}692 \]

\begin{center}
\begin{minipage}{0.9\textwidth}
\textbf{$B$ müsste in der ersten Runde zu $69$ Prozent wirken --- deutlich mehr als $A$ mit $60$.}
\end{minipage}
\end{center}

\schritt{Schritt 3 --- Warum die Hürde so hoch ist}
\begin{center}
\begin{minipage}{0.88\textwidth}
\kursiv{\textbf{Weil die Reihenfolge zweimal wirkt: $B$ müsste nicht nur besser wirken als $A$, sondern auch den schlechteren Rückfallweg ausgleichen.} Nach $B$-Versagen wirkt $A$ nur zu $30$ Prozent, nach $A$-Versagen wirkt $B$ zu $50$. Wer nur die erste Runde vergleicht, übersieht diesen zweiten Beitrag --- und das ist der häufigste Fehler beim Aufstellen von Therapiestufen.}
\end{minipage}
\end{center}
\end{aufgabe}

\begin{aufgabe}
\boxtitel{Aufgabe 7.2 --- Drei Runden}
Nach zweimaligem Versagen gibt es eine dritte Behandlung $C$ mit $20$ Prozent Erfolgsaussicht, Erfolg dann $0{,}80$ wert. Ändert das die Reihenfolge?

\loesung
\schritt{Schritt 1 --- Der Wert nach zwei Versagen}
\[ 0{,}20\cdot0{,}80 + 0{,}80\cdot0{,}30 = 0{,}160+0{,}240 = 0{,}400 \]
\kursiv{Statt $0{,}30$ wie bisher.}

\schritt{Schritt 2 --- Zuerst $A$}
\[ 0{,}60 + 0{,}40\bigl[0{,}50\cdot0{,}90 + 0{,}50\cdot0{,}400\bigr] = 0{,}600+0{,}40\cdot0{,}650 = \mathbf{0{,}8600} \]

\schritt{Schritt 3 --- Zuerst $B$}
\[ 0{,}40 + 0{,}60\bigl[0{,}30\cdot0{,}90 + 0{,}70\cdot0{,}400\bigr] = 0{,}400+0{,}60\cdot0{,}550 = \mathbf{0{,}7300} \]

\schritt{Schritt 4 --- Der Vergleich}
\begin{center}
{\small
\renewcommand{\arraystretch}{1.48}
\begin{tabular}{@{}P{2.4cm}P{2.2cm}P{2.2cm}P{1.8cm}@{}}
\toprule
& \textbf{zwei Runden} & \textbf{drei Runden} & \textbf{Zuwachs}\\
\midrule
zuerst $A$ & $0{,}8400$ & $0{,}8600$ & $+0{,}020$\\
zuerst $B$ & $0{,}6880$ & $0{,}7300$ & $+0{,}042$\\
\bottomrule
\end{tabular}}
\end{center}

\begin{center}
\begin{minipage}{0.9\textwidth}
\textbf{Die Reihenfolge bleibt, aber der Abstand schrumpft.}
\end{minipage}
\end{center}

\begin{center}
\begin{minipage}{0.88\textwidth}
\kursiv{\textbf{Je mehr Rückfalloptionen es gibt, desto weniger hängt vom ersten Schritt ab} --- ein allgemeines Muster, das man auch aus der Antibiotikawahl kennt: Wo Reserven existieren, wiegt der erste Fehlgriff weniger. Deshalb ist die Reihenfolge dort am wichtigsten, wo die zweite Chance fehlt.}
\end{minipage}
\end{center}
\end{aufgabe}

\begin{aufgabe}
\boxtitel{Aufgabe 7.3 --- Der Wert des Wartens}
Statt sofort zu behandeln, kann man einen Monat abwarten: Mit $30$ Prozent bessert sich der Zustand von selbst (Nutzen $1{,}00$), sonst behandelt man mit unverändertem Erfolg, aber der Nutzen sinkt wegen der Verzögerung um $0{,}05$.

\loesung
\schritt{Schritt 1 --- Sofort behandeln}
\[ 0{,}60\cdot1{,}00 + 0{,}40\cdot0{,}600 = 0{,}8400 \]

\schritt{Schritt 2 --- Warten}
\[ 0{,}30\cdot1{,}00 + 0{,}70\cdot\bigl[0{,}8400 - 0{,}05\bigr] = 0{,}300 + 0{,}70\cdot0{,}790 = 0{,}300+0{,}553 = \mathbf{0{,}8530} \]

\begin{center}
\begin{minipage}{0.9\textwidth}
\textbf{Warten ist besser --- um $0{,}013$.}
\end{minipage}
\end{center}

\schritt{Schritt 3 --- Der Kipppunkt}
Bei welcher Selbstheilungsrate $h$ sind beide gleich?
\[ h + (1-h)\cdot0{,}790 = 0{,}840 \quad\Longrightarrow\quad 0{,}210\,h = 0{,}050 \quad\Longrightarrow\quad h = 0{,}238 \]

\schritt{Schritt 4 --- Die Einordnung}
\begin{center}
\begin{minipage}{0.88\textwidth}
\kursiv{\textbf{Abwarten ist eine Handlung wie jede andere und gehört als eigener Ast in den Baum.} Dass es in vielen Modellen fehlt, verzerrt sie systematisch zugunsten des Eingreifens --- besonders bei Erkrankungen mit hoher Selbstheilungsrate. Der Wert des Wartens hat zudem einen zweiten Bestandteil, den dieses Beispiel unterschlägt: Man erfährt inzwischen mehr. Das ist der Informationswert aus Kapitel~3, und er macht das Warten oft noch günstiger.}
\end{minipage}
\end{center}
\end{aufgabe}

\begin{geschichte}
\boxtitelg{Ein Name gegen die Bürokratie}
\Person{Richard Bellman} entwickelte die dynamische Programmierung ab 1949 bei der RAND Corporation.

\begin{center}
\begin{minipage}{0.9\textwidth}
\kursiv{Den Namen wählte er nach eigener Auskunft, weil sein Vorgesetzter im Verteidigungsministerium das Wort \enquote{Forschung} nicht ertragen konnte --- \enquote{dynamisch} klang nach Tatkraft, \enquote{Programmierung} nach Planung. \textbf{Mit dem Rechnen am Computer hat das Wort nichts zu tun.} Bellman prägte auch den \enquote{Fluch der Dimension} aus Band 22, Kapitel~10 --- und zwar für genau dieses Verfahren: Die Zustandsmenge wächst mit jeder Größe exponentiell.}
\end{minipage}
\end{center}

\schritt{Und in der Medizin}
\begin{center}
\begin{minipage}{0.9\textwidth}
\kursiv{Die ersten medizinischen Anwendungen betrafen in den 1970er Jahren Screening-Intervalle und Transplantationswarteleisten. \textbf{Heute ist das Verfahren vor allem als Grundlage des bestärkenden Lernens bekannt} --- Band 22, Kapitel~15 --- und wird auf Behandlungsstrategien angewandt, wo es unmittelbar auf die Grenzen aus Band 23, Kapitel~12 trifft.}
\end{minipage}
\end{center}
\end{geschichte}

\begin{forschung}
\boxtitelp{Sequenzen heute}
\begin{center}
{\small
\renewcommand{\arraystretch}{1.45}
\begin{tabular}{@{}P{3.0cm}P{5.2cm}@{}}
\toprule
\textbf{Stufenpläne} & Therapiesequenzen als Strategie statt als Reihe von Einzelentscheidungen.\\
\textbf{SMART-Studien} & Versuchsaufbauten, die an jedem Schritt neu auslosen.\\
\textbf{Abwarten} & Der Wert des Zuwartens als eigener Forschungsgegenstand.\\
\bottomrule
\end{tabular}}
\end{center}

\schritt{Zur zweiten Zeile}
\begin{center}
\begin{minipage}{0.88\textwidth}
\kursiv{In einer sequenziell randomisierten Studie wird nach Versagen der ersten Therapie erneut ausgelost. \textbf{Damit lassen sich ganze Strategien vergleichen --- nicht nur einzelne Mittel} ---, und zwar unter den Bedingungen von Band 23, Kapitel~3: Das Los stellt an jedem Schritt her, was sonst behauptet werden müsste. In der Psychiatrie war die STAR*D-Studie zur Depressionsbehandlung ein früher großer Vertreter.}
\end{minipage}
\end{center}
\end{forschung}

\begin{merksatz}
\boxtitel{Was Sie jetzt können}
\begin{center}
\renewcommand{\arraystretch}{1.5}
\begin{tabular}{@{}P{3.6cm}P{5.0cm}@{}}
\toprule
\textbf{Satz 7.1} & Bellman: sofort plus Rest\\
\textbf{Optimalitäts\-prinzip} & jede Fortsetzung ist selbst optimal\\
\textbf{Ersparnis} & linear statt doppelt exponentiell\\
\textbf{Reihenfolge} & wirkt zweimal\\
\textbf{Warten} & ist eine Handlung\\
\bottomrule
\end{tabular}
\end{center}

\begin{center}
\begin{minipage}{0.88\textwidth}
\kursiv{Endlich viele Schritte lassen sich rückwärts rechnen. Das nächste Kapitel behandelt Entscheidungen ohne Ende --- und findet, dass dieselbe Gleichung dann eine Lösung hat, die man suchen muss statt sie aufzurollen.}
\end{minipage}
\end{center}
\end{merksatz}
